\documentclass[11pt]{article}
\usepackage[a4paper,margin=1in]{geometry}
\usepackage{fontspec}
\usepackage[boldfont=false]{xeCJK}
\setCJKmainfont{FandolSong-Regular.otf}
\usepackage{amsmath}
\usepackage{amssymb}
\usepackage{array}
\usepackage{booktabs}
\usepackage{graphicx}
\usepackage{adjustbox}
\usepackage{enumitem}
\setlist[itemize]{topsep=2pt,partopsep=0pt,parsep=0pt,itemsep=2pt,leftmargin=1.4em}
\setlist[enumerate]{topsep=2pt,partopsep=0pt,parsep=0pt,itemsep=2pt,leftmargin=1.6em}
\usepackage{parskip}
\usepackage{fvextra}
\fvset{breaklines=true,breakanywhere=true,fontsize=\small}
\setlength{\parindent}{0pt}

\begin{document}

\begin{center}
  {\LARGE\bfseries Government macroeconomic intervention (A Level)}\\[0.35em]
  {\large A Level Economics}
\end{center}
\vspace{0.6em}

This topic pulls the macroeconomics together: the goals a government has, how they clash, and how well its tools work.

\section*{Objectives and the tools to reach them}

Governments aim at several \textbf{macroeconomic objectives} 宏观经济目标 at once:

\begin{center}
\adjustbox{max width=\linewidth}{%
\begin{tabular}{ll}
\toprule
\textbf{Objective} & \textbf{How it is measured} \\
\midrule
\textbf{economic growth} 经济增长 & the \% change in real GDP \\
\textbf{price stability} 物价稳定 (low \textbf{inflation} 通货膨胀) & the \% change in the consumer price index \\
\textbf{full employment} 充分就业 & the unemployment rate \\
\textbf{balance of payments} 国际收支 stability & the current account balance \\
\bottomrule
\end{tabular}}
\end{center}

To reach these goals, the government uses three sets of tools:

\begin{itemize}
\item \textbf{fiscal policy} 财政政策 — changing government spending and taxes.

\item \textbf{monetary policy} 货币政策 — changing the \textbf{interest rate} 利率 and the money supply.

\item \textbf{supply-side policy} 供给侧政策 — raising the economy's capacity and productivity.

\end{itemize}

Fiscal and monetary policy mainly work on \textbf{aggregate demand} 总需求, so they are quick but can cause inflation. Supply-side policy works on capacity, so it is slow but lasting.

\section*{Links and trade-offs between objectives}

The objectives are \textbf{interrelated} — reaching one can move another. Often there is a \textbf{trade-off} 取舍, a \textbf{conflict} 冲突 where gaining on one goal costs you on another.

\begin{itemize}
\item \textbf{unemployment versus inflation.} Raising demand to cut unemployment tends to push prices up. The \textbf{Phillips curve} 菲利普斯曲线 shows this short-run trade-off. (In the long run it may disappear.)

\item \textbf{growth versus inflation.} Fast demand-led growth can overheat the economy and cause inflation.

\item \textbf{growth versus the balance of payments.} As incomes rise, people buy more imports, so the current account can worsen.

\item \textbf{growth versus the environment.} More output often means more pollution and use of resources.

\item \textbf{growth versus equality.} Growth does not always reach everyone, so the income gap can widen.

\end{itemize}

Because the goals are linked, a government must often accept "second best" on one aim to protect another. A clever policy mix tries to ease the trade-offs — for example, supply-side policy can raise growth \textbf{and} lower inflation, easing the usual clash.

\section*{The effectiveness of policy}

Each tool has strengths and weaknesses.

\subsection*{Fiscal policy}

It can act directly and powerfully on demand, and it can target particular groups or regions. But:

\begin{itemize}
\item it suffers \textbf{time lags} 时滞 — a budget is set once a year, and the effects take time.

\item large deficits add to the \textbf{national debt} 国债.

\item it can cause \textbf{crowding out} 挤出效应 — heavy government borrowing raises interest rates and reduces private investment.

\end{itemize}

The \textbf{Laffer curve} 拉弗曲线 is a useful warning. It shows that as the tax rate rises, \textbf{tax revenue} 税收收入 first rises but then \textbf{falls}: very high rates discourage work and encourage tax avoidance, so revenue drops. There is a tax rate that raises the most revenue, and going above it is self-defeating.

\subsection*{Monetary policy}

It is flexible — the interest rate can be changed at any time. But it also works with a long time lag, and in a deep slump even very low rates may fail to make worried firms and households borrow.

\subsection*{Supply-side policy}

It is the only tool that can raise growth and lower inflation together, with no trade-off. But it is slow (education and training take years), it is costly, and some measures (like cutting benefits) can widen inequality.

\subsection*{Putting it together}

No single policy can hit every target. The objectives conflict, the effects come with delays, and the future is uncertain. So governments usually combine \textbf{demand-side policies} 需求侧政策 (fiscal and monetary) for the short term with supply-side policy for the long term, and accept that some goals must be traded against others.


\end{document}
